% paper2_zh.tex — 中文版 · 编译: xelatex paper2_zh; bibtex paper2_zh; xelatex ×2
\documentclass[11pt]{article}
\usepackage[margin=2.6cm]{geometry}
% no xeCJK in this distribution: raw XeTeX CJK setup via fontspec
\usepackage{fontspec}
\setmainfont{Noto Serif CJK SC}
\setsansfont{Noto Sans CJK SC}
\XeTeXlinebreaklocale "zh"
\XeTeXlinebreakskip = 0pt plus 1pt minus 0.1pt
\usepackage{graphicx}
\usepackage[export]{adjustbox}
\usepackage{float}
\usepackage{amsmath,amssymb,bm}
\usepackage[colorlinks=true,linkcolor=blue,citecolor=blue,urlcolor=blue]{hyperref}
\newenvironment{acknowledgments}{\paragraph*{致谢}}{}
\renewcommand{\abstractname}{摘\quad 要}
\renewcommand{\refname}{参考文献}
\renewcommand{\figurename}{图}
\renewcommand{\tablename}{表}

\graphicspath{{figures/}}
\begin{document}

\title{六方氮化硼斐波那契同位素纳米带中的拓扑声子能隙}
\author{吴荣华\\[-1pt]\small 珠海九重天航空科技有限公司，珠海，中国\\[-2pt]\small\texttt{wrh@gdtzn.com}
\and 陈汝清\\[-1pt]\small GUT Geoservice Inc.，蒙特利尔，加拿大\\[-2pt]\small\texttt{ruqing@hotmail.com}}
\date{\today}
\maketitle

\begin{abstract}
超均匀同位素序通过一条对横向通道平均免疫的幂律压制声子散射
[R.~Wu 与 R.~Chen, Zenodo (2026)]。本文将目标反转，追求相反的谱学对象：
受拓扑标号保护的声子能隙——而它恰恰会被同一种平均抹平。能隙的存活因此
要求严格单模的纳米带，并随之直面边缘粗糙度退相干。我们把粗糙度散射锚定
于一套网格收敛的边界贴合格林函数计算（经解析 Born 结果验证），发现在单模
波动极限 $\lambda \gg W$ 下，Ziman 射线公式将相干长度低估达一个数量级：
零参数定律 $\xi_{\mathrm{edge}} = 2W^2/[q^2 S_w(2q)]$ 在全部三个标度指数
（频率指数至 $10^{-3}$）与 $3.5\%$ 的中位幅值上复现 15 点数值扫描（残差
为经独立确认的二阶修正），把边缘平整度要求从原子级（$\lesssim 0.13$~nm）放宽到
光刻级（$\lesssim 1$~nm）。周期 $24.8$~nm 的斐波那契 $^{10}$BN/$^{11}$BN
调制在 $30$~nm 宽的纳米带中把一条陈数标号的能隙置于 $0.250$~THz，相对周期
超晶格的局域化长度代价为 $2.20$ 倍。能隙标号以可证伪方式确认（残差中位数
$3\times10^{-6}$，低于随机对照 $1385$ 倍），相位子泵浦呈现完整的反向传播
边界态结构：全部十二条主隙（$|q|$ 最高至 $11$）的逐端谱流均为 $\pm q$。
我们提出基于异步光学采样的声学腔测量方案，并内置一条陈数反号的对照能隙。
\end{abstract}

\section{引言}

在同系列前文~\cite{paper1}中我们证明：空间关联的同位素排布给出声子弛豫率
$\tau^{-1}_{\rm iso} \propto \omega^{\,d+1+\alpha}$，其中排布的超均匀指数
$\alpha$~\cite{torquato2003} 以与维度无关的可加方式修正
Rayleigh--Tamura 指数~\cite{tamura1983}。那篇文章的结论是压制性的：给同位素
排序会\emph{移除}长波散射。本文把目标反转，随之反转了由其导出的每一条几何
结论。

\textbf{镜像。}幂律是标度不变量：对三维样品的横向通道求平均只改前因子、
不动指数——这份免疫为文献~\cite{paper1}的厚样品几何背书，并使纳米带在
那里不仅无用而且有害。本文要的是能隙。能隙不是标度不变量：它在横向通道
$m$ 中的位置由该通道自己的纵向波矢 $q_m^2 = (\omega^2 - v^2 k_m^2)/v^2$
决定，通道求和逐通道迁移隙位，把康托集能隙抹成无特征的连续谱。保护幂律
的那种平均，正是摧毁能隙的那种平均。横向量子化因此不是偏好而是前提，且
必须取其最严格形式：单模运行，$N_{\rm prop} = 1$。这只允许一维纳米带——
正是文献~\cite{paper1}抛弃的几何——并把它绕开的边缘退相干问题交到我们
手里。第~\ref{sec:survival}节即为这场对峙。两篇文章并不矛盾：它们报告的
是同一件物理事实——通道平均保指数、灭能隙——只是读法对应于相反的目标。

\textbf{为何是斐波那契，而不是素数。}能隙本身并不拓扑。一维准周期链的
拓扑内容依赖能隙标号定理~\cite{bellissard1992}，它适用于具有纯点衍射谱与
切割投影构造的序列：斐波那契~\cite{kohmoto1983,ostlund1983}、Sturm 词、
二次无理数。在切割投影下，链是二维周期母晶格的一个切片，切片偏移 $\phi$
是一个合成的第二维度~\cite{kraus2012}，每条能隙携带整数标号 $q$，即母晶格
的陈数。素数不具备任何这类结构：其衍射更接近极限周期而非准周期，没有
相位子、没有合成维度、没有能隙标号。以素数为指标的调制会打开一族稠密的
弱能隙，却不携带任何不变量。我们明说这一点，因为这类构造的算术魅力可观，
而其拓扑内容为零。

若斐波那契同位素纳米带的能隙确受标号保护，其位置便由序列本身而非任何
无序实现所固定，且边界分支必须在一个相位子周期内恰好穿越 $q$ 号隙 $q$
次~\cite{kraus2012,apigo2019}。第~\ref{sec:topo}节检验这两条陈述；
第~\ref{sec:survival}节先确立这样的纳米带能否存在。

\section{生存窗口}
\label{sec:survival}

\subsection{转移矩阵与局域化长度}

质量为 $M_n$、近邻力常数为 $K$、间距为 $a$ 的链满足
\begin{equation}
\begin{pmatrix} u_{n+1} \\ u_n \end{pmatrix} = T_n
\begin{pmatrix} u_n \\ u_{n-1} \end{pmatrix}, \quad
T_n = \begin{pmatrix} 2 - M_n\omega^2/K & -1 \\ 1 & 0 \end{pmatrix},
\end{equation}
逆局域化长度即李雅普诺夫指数
$\gamma(\omega) = \lim_{N\to\infty} N^{-1}\ln\|\prod_n T_n\|$，
$\xi = a/\gamma$，隙内透射
$\mathcal{T}(L) \sim e^{-2L/\xi}$~\cite{anderson1958}。在
Born--Oppenheimer 近似下力常数与同位素无关，故单胞计算一次即服务所有
排布；只需重建对角质量矩阵。

\subsection{由质量对比度定隙强}

$^{10}$BN 与 $^{11}$BN 单元（24.01 对 25.01~amu）给出密度对比
$\Delta\rho/\rho = 4.1\%$、声阻抗对比 $\Delta Z/Z = 2.04\%$、单界面反射
$|r| = 0.0102$~\cite{vuong2018}。弱反射理论预言
\begin{equation}
\xi_{\rm gap} = \frac{d}{2|r|}, \qquad
\frac{\Delta\omega}{\omega} = \frac{2|r|}{\pi} = 0.67\%,
\end{equation}
$d$ 为调制周期；$2\times10^5$ 格点链上的直接李雅普诺夫计算把二者复现到
$1$--$6\%$。

\subsection{拓扑的代价}

准周期性把反射强度分摊到康托集能隙上，而非集中于一条 Bragg
共振~\cite{kohmoto1983,merlin1985}。其价格由同频对比裁定：Bragg 隙位于
$0.25$~THz 的周期超晶格（$d_{\rm per} = 40.1$~nm）给出
$\xi_{\rm per} = 1.96\,\mu$m，与解析式 $d/2|r|$ 吻合到三位有效数字；置于
同一频率的斐波那契 $(0,1)$ 隙给出 $\xi_{\rm gap} = 4.31\,\mu$m。在匹配的
频率与质量对比度下，拓扑的代价为
\begin{equation}
\xi_{\rm fib}/\xi_{\rm per} = 2.20 .
\end{equation}

\subsection{两种微扰，相反的对称性}

同位素调制横贯整个纳米带截面，在横向模基下严格对角：
$V^{\rm iso}_{mm'} = -\omega^2\epsilon(z)\delta_{mm'}$——它在每个通道里
独立开隙、隙位随通道移动，即前述抹平机制的通道语言重述。边缘粗糙度破坏
横向平移，是非对角的；其后果在性质上取决于模数。多模运行时它把声子散射
到其他子带——真正的退相干。单模运行时无处可散，只剩同通道背散射
$q \to -q$，机制退化为一维 Anderson 局域化：它\emph{加}在 $\gamma$ 上，
而非摧毁能隙。单模纳米带因此不只是量变的更安全，而是把主导损耗换成了
另一种性质的东西；确立严格单模运行由此成为本节的核心技术要求。

\subsection{张力与工作点}
\label{sec:tension}

自由（Neumann）边缘的单模条件为 $f < v/2W$；相干要求
$\xi_{\rm dis} > 3\,\xi_{\rm gap}$。以 $\xi_{\rm gap} = C_\xi\lambda$ 与
第~\ref{sec:calib}节的标定定律合并二者，得
\begin{equation}
\lambda^3 > 24\pi^2 C_\xi\, S_w(2q) \;\Rightarrow\;
\lambda \gtrsim 54\,[\Delta_r/\mathrm{nm}]^{2/3}~\mathrm{nm},
\end{equation}
远软于 Ziman 形式所隐含的线性边界；从 $\Delta_r$ 到 $\Delta_r^{2/3}$ 的
指数变化正是下文 $W^2$ 定律的直接推论。

斐波那契链的能隙落在 Fourier 模上：
\begin{equation}
f_{(p,q)} = \frac{v}{2d}\,\bigl|p + q\,\tau^{-1}\bigr|,
\label{eq:module}
\end{equation}
$\tau = (1+\sqrt5)/2$。令 $q = 1$ 隙落在目标频率即定
$d = \tau^{-1} v/2f = 24.8$~nm（每字母 $99$ 个格点）；模导向扫描确认四条
主隙全部落在式~\eqref{eq:module}预言位置的 $0.3\%$ 以内，标号经积分态
密度（IDOS）核实。两个发现值得记录。其一，$(1,0)$ 隙——周期直觉唯一会
去找的那条，位于 $v/2d$——在此装饰下比 $(0,1)$ 弱三十倍，且携带
$q = 0$：不拓扑，而它几近缺席实属幸事。其二，$0.154$~THz 处的 $(1,-1)$
隙与 $(0,1)$ 同强（$\xi = 4.18\,\mu$m）而标号相反；二者均在 $30$~nm
纳米带的单模顶棚之下，故单个器件同时承载两条陈数相反的拓扑能隙——为
第~\ref{sec:expt}节提供了内置对照。

\begin{table}[H]
\centering
\caption{工作点。$\xi_{\rm dis}$ 取自第~\ref{sec:calib}节的标定定律；
余量为 $\xi_{\rm dis}/\xi_{\rm gap}$。}
\begin{tabular}{lc}
\hline\hline
频率 / 工作隙 & $0.250$~THz，$(0,1)$，$q = 1$\\
带宽 $W$ / 周期 $d$ & $30$~nm / $24.8$~nm\\
$\xi_{\rm gap}$ & $4.31\,\mu$m\\
$\Delta_r = 0.6$ ($1.0$)~nm 时的 $\xi_{\rm dis}$ & $76$ ($30$)\,$\mu$m\\
相干余量 & $18$ 倍（$7.0$ 倍）\\
\hline\hline
\end{tabular}
\end{table}

\subsection{标定：射线公式的失效}
\label{sec:calib}

求解器在任何扫描之前经三重独立验证：精确复现干净带极限（三个网格上
$\mathcal{T} = N_{\rm prop}$ 至五位，每个透射本征值均为 1）；通过实现匹配
的网格收敛门（$h = 1$ 与 $0.5$~nm 之间 $1.3\%$）；并复现解析单凸起 Born
反射 $|r|^2 = (q^2/W^2)|\tilde a(2q)|^2$，随网格收敛到均匀的 $+3.7\%$
超额——一个真实的二阶修正，其作用在下文起到判决性作用。

定版扫描——十五个单模锁定参数点、每点 $24$ 个无序种子、六个长度，且每个
长度窗口横跨定律自身 $\xi$ 的 $(0.4$--$3)$ 倍，从而真正采样衰减而非外推
——给出
\begin{align}
\frac{{\rm d}\ln\xi}{{\rm d}\ln\Delta_r} &= -1.91 \pm 0.06, &
\frac{{\rm d}\ln\xi}{{\rm d}\ln W} &= +1.87 \pm 0.08,\nonumber\\
\frac{{\rm d}\ln\xi}{{\rm d}\ln f} &= -1.384 \pm 0.095,
\end{align}
对照定律的 $(-2,\,+2,\,-2+4q^2L_c^2 = -1.38)$：三者全部落入容差，频率
指数命中到千分之一。绝对值以中位 $\xi_{\rm meas}/\xi_{\rm law} = 0.965$
吻合，而单凸起测试对同一二阶超额独立预言 $1/1.037 = 0.964$：对同一个修正
的两次测量——一次确定性、一次系综平均——吻合到 $10^{-3}$。纵观全扫描，
Ziman 公式把 $\xi$ 低估达中位 $10.8$ 倍：在单模波动极限 $\lambda \gg W$
下，粗糙边缘是 Born 透明的，壁面漫反射碰撞的图像并不适用。其标度结构——速率 $q^2 S_w(2q)/W^2$——与
Santamore 与 Cross~\cite{santamore2001} 在热导量子背景下对少模弹性波导
粗糙度散射的微扰处理一致（彼处为夹紧弹性边界）；式~\eqref{eq:law}的零
参数闭式、其绝对归一以及深入局域化区的验证，是本文新增的部分。残差保留一个
随 $qL_c$ 的温和趋势（从 $qL_c = 0.16$ 处的 $0.78$ 到 $0.42$ 处的
$1.12$）——正是超出首阶定律的二阶结构应有的指纹。

方法学记录——三种失败模式，各由不同的门捕获，无一能凭目检发现。阶梯掩模
求解器产生随网格发散的伪\emph{过量}散射（$h = 2/1/0.5$~nm 下
$\xi = 120/42/16$~nm，无连续极限），由网格收敛检验捕获。变换算符的朴素
有限差分离散化非厄米且破坏流守恒（透射本征值高至 $1.26$），却在网格上
\emph{收敛}且总透射与 Born 相符：网格收敛检验对它失明，唯有作用于透射
本征值谱的幺正门将其揭露；治法是从问题的二次型出发装配离散算符，使厄米性
（$R_{m+1} = P_m^{\mathsf T}$ 按构造成立）与精确的斜壁 Neumann 条件成为
构造本身的性质而非手推格式的性质。长度窗口若系于 Ziman 估计——短了十一
倍——则只采样浅衰减，曾在初测中把 $W$ 指数偏移 $+0.3$：窗口必须追踪被
检验的定律本身。另有一次梯子跨越横模阈值的扫描在同一轴上返回符号相反的
指数（$+5.45$ 与 $-2.91$），即单模条件未受强制的诊断指纹。

其后果是本节的核心实用结论：相干性不再要求 $\Delta_r \lesssim 0.13$~nm
——单原子行——而是在 $\Delta_r = 0.6$~nm 下仍余十八倍裕度、在
$1.0$~nm 下越过三倍线（表~1、图~\ref{fig:survival}）：落入常规电子束
光刻的范围。

在第~\ref{sec:topo}节的拓扑建立于这些数字之上以前，有一段适用范围声明
须先交代给读者：以上全部处理的是速度 $v = 20$~nm/ps 线性色散的单一标量
极化。真实的悬空 h-BN 纳米带携带四支低频声子——纵向、面内横向、扭转，
以及色散为二次、主导低频态密度的弯曲支。本文依赖的两个机制是极化无关
的：能隙标号只要求质量序列的一维准周期性，Born 透明性只要求弱边界微扰
的 $q$ 选择性背散射，故式~\eqref{eq:law}与模~\eqref{eq:module}的结构
可逐支迁移。数字则不然：支间模式混合、弯曲支的二次色散与逐支不同的边缘
耦合将移动工作点频率、$\xi$ 值与余量；标量模型最恰当的读法，是纵向
通道的定量处理兼其余各支的结构模板。完整的弹性计算是自然的下一步，
超出本文范围。

\begin{figure}[H]
\includegraphics[width=0.8\textwidth,center]{fig_survival.pdf}
\caption{\label{fig:survival}生存相图，全部解析自式~\eqref{eq:law}：
单模饱和带宽下 $(f, \Delta_r)$ 平面内的余量
$\xi_{\rm dis}/3\xi_{\rm gap}$，等值线 $M = 1, 3, 10$。灰色点划线为
Ziman 射线公式将强加的边界。星形：工作点（$\Delta_r = 0.6$~nm，
$21$ 倍）；圆点：典型 EBL 边缘（$1.0$~nm，$7.4$ 倍）。}
\end{figure}

\section{拓扑验证}
\label{sec:topo}

\subsection{可证伪的能隙标号}

隙内 IDOS 为常数，且被能隙标号定理限制为
$\mathcal{N} = (p + q\,\tau^{-1}) \bmod 1$~\cite{bellissard1992}。我们以
转移矩阵递推上的 Sturm 节点计数求 $\mathcal{N}$（$\mathcal{O}(N)$，无需
对角化）；表~2 列出 $6\times10^4$ 格点链的十二条最强隙
（图~\ref{fig:gamma}）。最强的两条取教科书值 $\tau^{-1}$ 与
$1 - \tau^{-1}$。

若照字面理解，该检验是空洞的：$\mathbb{Z} + \tau^{-1}\mathbb{Z}$ 在
$[0,1]$ 中稠密，任何数都有某对整数拟合之。检验的内容是定量的。以完全
相同的拟合流程（$|q| \le 30$）处理 $5\times10^3$ 个均匀随机对照：实测隙
残差中位数 $2.95\times10^{-6}$（最大 $5.95\times10^{-6}$），随机对照
$4.08\times10^{-3}$——分离 $1385$ 倍；随机值拟合得不差于\emph{最差}
实测隙的概率为 $p = 1.4\times10^{-3}$。我们建议凡数值宣称能隙标号者皆
报告此对照。

\begin{figure}[H]
\includegraphics[width=0.8\textwidth,center]{fig_gamma_idos.pdf}
\caption{\label{fig:gamma}斐波那契链的李雅普诺夫谱 $\gamma(\omega)$
（上）与 IDOS 阶梯（下），标出表~2 的十二条标号隙。插图：标号残差，
实测隙对 $3\times10^3$ 随机对照（对数轴）——本节可证伪性论证的单图
呈现。}
\end{figure}

\begin{table}[H]
\centering
\caption{十二条主隙（$N = 6\times10^4$ 格点）。$\xi$ 取隙心；残差为
$|q| \le 30$ 下最优对的 $|\mathcal{N} - (p + q\tau^{-1})|$。}
\begin{tabular}{cccccc}
\hline\hline
$f$ (THz) & $\gamma$ (每格点) & $\xi$ (nm) & IDOS & $(p,q)$ & 残差\\
\hline
14.386 & 0.00839 & 29.9 & 0.38197 & $(1,-1)$ & $7{\times}10^{-7}$\\
17.213 & 0.00290 & 86.5 & 0.47213 & $(-2,4)$ & $3{\times}10^{-6}$\\
18.792 & 0.00356 & 70.3 & 0.52787 & $(3,-4)$ & $3{\times}10^{-6}$\\
19.693 & 0.00170 & 147.4 & 0.56230 & $(-5,9)$ & $6{\times}10^{-6}$\\
21.040 & 0.01744 & 14.4 & 0.61803 & $(0,1)$ & $7{\times}10^{-7}$\\
22.222 & 0.00284 & 88.1 & 0.67377 & $(5,-7)$ & $5{\times}10^{-6}$\\
22.859 & 0.00345 & 72.6 & 0.70820 & $(-3,6)$ & $4{\times}10^{-6}$\\
23.760 & 0.01155 & 21.7 & 0.76393 & $(2,-2)$ & $1{\times}10^{-6}$\\
24.222 & 0.00185 & 135.1 & 0.79837 & $(-6,11)$ & $7{\times}10^{-6}$\\
24.479 & 0.00239 & 104.7 & 0.81967 & $(7,-10)$ & $7{\times}10^{-6}$\\
24.826 & 0.00812 & 30.8 & 0.85410 & $(-1,3)$ & $2{\times}10^{-6}$\\
25.239 & 0.00516 & 48.5 & 0.90983 & $(4,-5)$ & $3{\times}10^{-6}$\\
\hline\hline
\end{tabular}
\end{table}

\subsection{相位子泵浦：完整的反向传播结构}
\label{sec:pump}

切割投影词
$s_n(\phi) = \lfloor (n{+}2)\tau^{-1}{+}\phi \rfloor -
\lfloor (n{+}1)\tau^{-1}{+}\phi \rfloor$ 的相位子角 $\phi$ 即合成维度；
完整一周遍历整个包络，边界分支必须恰好穿越 $q$ 号隙 $q$
次~\cite{kraus2012}。追踪此类分支需要谨慎：按频率排序会在能带交叉处
交换态的身份，故我们以本征矢重叠配合匈牙利指派延拓态的身份，最优重叠
低于 $0.15$ 的分支即终止而非强行匹配，并分别记录左、右端权重
（图~\ref{fig:phason}）。

对 $(0,1)$ 隙的首轮追踪（$610$ 格点自由链、$400$ 步）在窗口内追出 $27$
条分支，报告仅一条边界束缚分支——右端、权重 $0.721$、扫过隙宽 $86\%$、
一处拐折——且在链长 $610$--$1597$、步数 $400$--$800$ 间逐位稳定。单看
此结果会得出单侧泵浦的结论。两项与基无关的测量纠正了这一读法。
\emph{谱流}：在 $600$ 个相位子步上数各隙心以下的本征值，十二条隙每一条
都在一周内恰好出现 $|q|$ 次下行与 $|q|$ 次上行穿越——$(0,1)$ 为
$1{+}1$，$(-1,3)$ 为 $3{+}3$，直至 $(-6,11)$ 的 $11{+}11$——无一例外。
\emph{穿越时刻分类}：在每个探测到的穿越处携本征矢对角化并直接读取边缘
权重，所有 $q > 0$ 隙的下行穿越均右端束缚、上行均左端束缚，$q < 0$ 严格
镜像；$(-1,3)$ 的六次穿越完美交替换边，$(0,1)$ 对偶对的权重为 $0.72$ 与
$0.71$。反向传播对完整存在；逐端谱流等于 $\pm q$；体--边对应以标准形式
对全部十二个标号成立。

\begin{figure}[H]
\includegraphics[width=0.8\textwidth,center]{fig_phason.pdf}
\caption{\label{fig:phason}$(0,1)$ 隙的相位子谱：频率对 $\phi$，体态为
散点，边界束缚分支按\emph{瞬时}主导端着色（蓝为左端、橙为右端），线宽
$\propto$ 边缘权重。隙内简并处的颜色切换即第~\ref{sec:topo}C 小节的
身份互换。}
\end{figure}

\subsection{关于重叠追踪的告诫}

首轮为何只见一条分支？两支反向传播边界模在隙内简并处杂化为携带
\emph{双端}权重的对称与反对称组合；在该避免交叉处，重叠延拓——依其
构造——跟随光滑形变的杂化态而交换分支身份，把两支的下降段缝成一条带
单个拐折的表观分支：恰是首轮报告的全部现象。第~\ref{sec:pump}小节的
谱流分类从不要求相邻步之间的连续性，故对此免疫。固定（Dirichlet）端下
对偶对依然存在（穿越位于 $\phi = 0.833, 0.932$，自由端为
$0.618, 0.147$）：终端条件只挪动穿越位置，未移除任何东西。我们记录此
伪象，因为匈牙利重叠延拓是泵浦谱分析的标准做法，而它恰在拓扑泵浦所
必然制造的场景——反向传播分支于隙内相遇——中静默失效。

\section{实验设计}
\label{sec:expt}

\subsection{频段选择；瞬态热栅之排除}

工作隙位于 $0.250$~THz，$\lambda = 80$~nm。亚太赫兹频段天然的波矢选择性
探针——瞬态热栅——在此被光学而非声学排除：波长 $\lambda_{\rm opt}$ 的
双光束无法写出细于 $\lambda_{\rm opt}/2$ 的光栅，$400$~nm 泵浦封顶于约
$200$~nm——粗了 $2.5$ 倍——三倍频 $266$~nm 也只达 $133$~nm。任何几何
都无法关闭这个因子。测量必须在时域进行：两台重复率相差数 kHz 的飞秒
振荡器之间的异步光学采样（ASOPS）~\cite{bartels2007}在完整的 $12.5$~ns
重复窗口内提供约 $100$~fs 的有效分辨率，将 $4$~ps 的声学周期采样约
$3\times10^3$ 个循环——谱分辨率约 $0.3$~GHz。

\subsection{能隙的腔模谱学}

纳米带两端约 $20$~nm 的铝垫同时充当宽带声子源~\cite{thomsen1986}、探测
反射镜与声学法布里--珀罗腔的反射端；带体内部仅通过其模式谱受审——正如
共振探测诊断无法就地触及的结构。$L = 20\,\mu$m 的腔自由谱程为
$v/2L = 0.5$~GHz；$(0,1)$ 隙（$1.7$~GHz）驱逐三到四个连续腔模，且以
$\xi_{\rm gap} = 4.31\,\mu$m 计，隙内透射 $e^{-2L/\xi} \approx 10^{-4}$：
模式是被移除而非被移位。观测量取差分，与文献~\cite{paper1}相同：
斐波那契带与同芯片、同换能器、同同位素配比之随机排列对照带成对制备；
边界、换能器与非谐贡献一阶抵消（图~\ref{fig:expt}）。分辨率预算如实
陈述：$(0,1)$ 隙五个分辨单元、$(1,-1)$ 隙三个——足以确立模式缺失，
不足以做线型分析。

\begin{figure}[H]
\includegraphics[width=0.72\textwidth,center]{fig_experiment.pdf}
\caption{\label{fig:expt}测量方案（示意）。(a)~带斐波那契同位素字母与铝
换能器的纳米带。(b)~示意性腔谱：配比匹配的随机对照（灰，完整模梳）对
斐波那契（蓝，隙被清空），隙内为随 $\phi$ 移动的边界模。}
\end{figure}

\subsection{相位子阵列的边界态泵浦}

相位子角由光刻固定，一带一值；$\phi = 0, 0.05, \ldots, 0.95$ 的约 $20$
根带阵列把泵浦周期转换为芯片位置。第~\ref{sec:pump}小节预言一条隙内
共振，其频率沿阵列单调推进、在大半个周期内可见。两项结构性对照不请
自来。同一批带中的 $(1,-1)$ 隙携带相反标号，其隙内共振必须沿同一阵列
反向绕行；任何平凡系统误差——应变、厚度漂移、成分梯度——都同向移动
两条隙，从而被排除。反向传播结构亦可直接检验：下行共振（$(0,1)$ 隙的
$\phi \approx 0.62$）居于一端、上行共振（$\phi \approx 0.15$）居于
另一端，故互换泵浦与探测换能器必须交换阵列中隙内响应更强的区段——
单侧机制无法复现的指纹。

\subsection{样品规格}

$W = 30$~nm（单模至 $0.267$~THz）；$L = 20\,\mu$m
（$4.6\,\xi_{\rm gap}$）；$d = 24.8$~nm（每字母 $99$ 格点）；每带约
$800$ 个字母（$F_{16} = 987$ 词截断）；$\Delta_r \le 0.6$~nm rms
（$18$ 倍余量）；$T \le 4$~K。同位素图形为一维、面内，处于电子束掩模
分辨率之内，候选路线为双前驱体顺序再生长。与文献~\cite{paper1}的对比
值得重申：彼处设计要求 $300\,\mu$m 上 $2.4\times10^5$ 个独立指定单元，
超出现有一切技术；此处仅需 $20\,\mu$m 上约 $800$ 个介观尺度字母，且
边缘容差因第~\ref{sec:calib}节的物理放宽一个数量级。把工作点压向低频
的那份张力，同时把它送进了可制备的疆域。

\begin{acknowledgments}
数值计算全部采用开源软件完成。复现本文每个数字所需的全部代码、标定数据
与出图脚本见
\url{https://github.com/Ruqing1963/hbn-fibonacci-topological-phonons}。
\end{acknowledgments}

\bibliographystyle{unsrt}
\bibliography{refs}

\end{document}
